\appendix

\section{How the Training Data Was Generated}
\label{app:data-generation}

This appendix explains the data-generation pipeline used for the model in the
main paper. The goal of the generator is simple: create many small process
behaviors, write each behavior in several forms, and keep an explicit record of
which forms describe the same behavior. The model can then learn from paired
examples instead of guessing which log, tree, and Petri net belong together.

\subsection{What One Training Row Contains}
\label{app:data-row}

Each row in the generated data is a \emph{process sample}. A row contains four
main fields:
\begin{enumerate}
  \item a canonical process tree $T$;
  \item a generated log view $L$, stored as a finite list of visible activity
        traces;
  \item a Petri-net graph $G$, stored as typed nodes, labeled transitions,
        directed arcs, and initial/final marking vectors;
  \item an equivalence identifier $e$, shared by all rows that represent the
        same behavior.
\end{enumerate}
The process tree is the decoder target. The generated log is the trace-based
view of the same behavior. The Petri-net graph is the graph-based view. During
training, rows with the same equivalence identifier are treated as positive
matches for the contrastive objective.

The persisted JSONL rows also store bookkeeping fields:
\texttt{model\_variant\_id}, \texttt{log\_view\_id},
\texttt{representation\_kind}, \texttt{equivalence\_level}, and
\texttt{metadata}. These fields make it possible to inspect which Petri-net
variant was used, which generator motif produced the family, and which
equivalence checks were passed.

\subsection{Current Split and Generator Settings}
\label{app:data-settings}

The current data directory was generated with the behavior-family generator,
fixed seed 13, and two representation rows per behavior family. The split sizes
are:
\begin{center}
\small
\begin{tabular}{lrr}
\toprule
Split & Rows & Behavior families \\
\midrule
Training & 8{,}192 & 4{,}096 \\
Validation & 1{,}024 & 512 \\
Test & 1{,}024 & 512 \\
\bottomrule
\end{tabular}
\end{center}

The important generator settings are:
\begin{itemize}
  \item maximum ordinary-tree recursion depth: 8;
  \item minimum ordinary-tree depth/size targets: depth 4 and 20 nodes;
  \item maximum activity labels: 30;
  \item maximum operator arity: 3;
  \item traces per generated log view: 128;
  \item activity-label reuse probability: 0.15;
  \item early leaf probability during tree recursion: 0.65;
  \item representation rows per behavior family: 2;
  \item log views per behavior family: 1;
  \item controlled-motif context suffix: 4 fresh activities;
  \item strict minimum motif coverage: 8 families per motif in training and
        4 per motif in validation and test;
  \item enabled process-tree operators: \texttt{SEQ}, \texttt{XOR},
        \texttt{AND}, and two-child \texttt{LOOP}.
\end{itemize}
The \texttt{LOOP} operator is generated in the ordinary-tree family. The
generated and decoded loop form is the two-child form
$\operatorname{LOOP}(T_{\mathit{body}},T_{\mathit{redo}})$; the tokenizer mask
therefore emits \texttt{ARITY\_2} after a \texttt{LOOP} token. Because a loop
can repeat an arbitrary number of times, ordinary loop families use finite
process-tree playout to construct the observed trace pool rather than trying to
enumerate an unbounded language exhaustively.

\subsection{Why the Generator Uses Behavior Families}
\label{app:behavior-families}

The generator does not create isolated triples one by one. Instead, it creates
\emph{behavior families}. A behavior family is one underlying behavior plus two
Petri-net representations of that behavior. Both rows in a family share:
\begin{itemize}
  \item the same canonical process tree;
  \item the same generated trace list;
  \item the same equivalence identifier.
\end{itemize}
They differ in their Petri-net graph. Sometimes the second graph is only a
renamed copy of the first graph. In other cases, the second graph uses a
different internal structure while keeping the same complete visible traces.
This is important because the model should learn behavior, not only superficial
graph shape.

The current data is balanced across four family types. Each split contains the
same number of behavior families from each type:
\begin{center}
\small
\begin{tabular}{lrrr}
\toprule
Family type & Training families & Validation families & Test families \\
\midrule
Ordinary tree & 1{,}024 & 128 & 128 \\
Duplicate label vs. silent routing & 1{,}024 & 128 & 128 \\
Concurrency vs. explicit interleaving & 1{,}024 & 128 & 128 \\
Non-free-choice M pattern & 1{,}024 & 128 & 128 \\
\bottomrule
\end{tabular}
\end{center}

\subsection{Step 1: Choose the Family Type}
\label{app:choose-family-type}

For each split, the generator first computes how many behavior families are
needed. Because the current configuration stores two rows per behavior family,
8{,}192 training rows become 4{,}096 training families. The same calculation
turns 1{,}024 validation rows and 1{,}024 test rows into 512 families each.

The four family types have equal weight. The generator allocates a deterministic
quota to each type and shuffles the order using a seed derived from the global
seed, the split name, and the phrase \texttt{motif-quota-plan}. This gives a
mixed order of samples on disk while preserving exact class balance.

The current run uses strict class coverage. Before replacing the data
directory, the generator verifies that every requested row count is divisible
by the number of rows per family and that the requested family count can meet
the per-motif minimum. An infeasible request therefore fails before the
existing split files are removed. A separate \texttt{best\_effort} mode is
available for small diagnostic datasets; it records motif and representation-
slot deficits instead of requiring them to be zero.

A complete behavior family is assigned to exactly one split, and its
equivalence identifier includes that split. The generator does not, however,
deduplicate complete visible languages across different families or splits.
Consequently, family members cannot leak across a split boundary, but equal or
near-equal behaviors can still occur independently in training and evaluation.

\subsection{Step 2: Build the Canonical Process Tree}
\label{app:build-tree}

The canonical process tree is the target representation used by the decoder.
There are two cases.

\paragraph{Ordinary-tree families.}
For ordinary-tree families, the generator samples a random process tree. It
first samples a pool of activity labels, then recursively builds the tree:
\begin{enumerate}
  \item If the maximum depth has been reached, it creates an activity leaf.
  \item Otherwise, with probability 0.65 it also stops early and creates an
        activity leaf.
  \item If it does not stop, it chooses one operator from \texttt{SEQ},
        \texttt{XOR}, \texttt{AND}, and \texttt{LOOP}.
  \item For \texttt{SEQ}, \texttt{XOR}, and \texttt{AND}, it samples an arity
        between 2 and 3 and recursively generates the children.
  \item For \texttt{LOOP}, it recursively generates exactly two children: the
        body and the redo branch.
\end{enumerate}
Activity labels are introduced as \texttt{a0}, \texttt{a1}, and so on. With
probability 0.15 the generator reuses an earlier label instead of introducing a
new one. After the tree is built, activity labels are canonicalized to
\texttt{A0}, \texttt{A1}, and so on in first-seen order. Children of
\texttt{XOR} and \texttt{AND} nodes are sorted by a canonical structural key, so
equivalent child permutations do not become different targets.

If a generated ordinary tree has fewer than two unique activity labels, the
generator tries to make it less trivial by adding another activity in sequence.
The sampler retries ordinary trees and keeps the first tree that contains a
loop, reaches the configured minimum depth, and has at least 20 nodes. If this
does not happen within the retry budget, it keeps the deepest and largest
candidate encountered.

\paragraph{Controlled behavior families.}
The other three family types use small hand-defined canonical trees:
\begin{itemize}
  \item duplicate-label vs. silent-routing families use
        $\operatorname{XOR}(\operatorname{SEQ}(A0,A1),
        \operatorname{SEQ}(A0,A2))$;
  \item concurrency vs. explicit-interleaving families use
        $\operatorname{AND}(A0,A1)$;
  \item non-free-choice M-pattern families use
        $\operatorname{XOR}(A0,\operatorname{AND}(A1,A2))$.
\end{itemize}
These trees are intentionally small because their purpose is to isolate one
specific representation issue at a time. To avoid making every controlled case
too bare, the generator appends up to four fresh activity labels as a shared
sequential suffix. The suffix is added as nested binary sequence nodes, so the
configured maximum arity remains 3. For example, if a motif already uses
\texttt{A0}, \texttt{A1}, and \texttt{A2}, the suffix can add \texttt{A3}
through \texttt{A6}. The same suffix is added to both Petri-net variants in the
family.

\subsection{Step 3: Build Two Petri-Net Variants}
\label{app:build-petri-variants}

Each behavior family produces two Petri-net variants. These variants are the
reason the dataset has two rows per behavior family.

\paragraph{Ordinary-tree families.}
The first Petri net is obtained by converting the process tree with PM4Py. The
second Petri net is an isomorphic copy: the places and transitions are renamed,
but the arcs, labels, initial marking, and final marking preserve the same
behavior. This teaches the model not to treat arbitrary node names as semantic
information.

\paragraph{Duplicate-label vs. silent-routing families.}
Both variants accept the same two visible traces, $A0,A1$ and $A0,A2$ before the
shared suffix. One variant has two visible \texttt{A0} transitions, one for each
branch. The other variant has a single visible \texttt{A0} transition followed
by silent routing to the \texttt{A1} or \texttt{A2} branch. The visible behavior
is the same, but the graph explains the choice differently.

\paragraph{Concurrency vs. explicit-interleaving families.}
Both variants accept $A0,A1$ and $A1,A0$ before the shared suffix. The first
variant is the Petri net obtained from the parallel process tree, so the two
activities can run concurrently. The second variant lists the two possible
orders explicitly as two serial paths. The visible traces are the same, but one
graph contains true concurrency and the other encodes only the allowed
interleavings.

\paragraph{Non-free-choice M-pattern families.}
Both variants accept the same visible traces for
$\operatorname{XOR}(A0,\operatorname{AND}(A1,A2))$ before the shared suffix. The
first variant is the standard PM4Py conversion of the block-structured tree. The
second variant is a hand-built Petri net with a non-free-choice M pattern. This
gives the model examples where behavior is preserved even though the Petri-net
structure is less block-like.

\subsection{Step 4: Check Behavioral Equivalence}
\label{app:equivalence-check}

Before a non-loop behavior family is accepted, the generator compares the
visible trace sets returned for all of its Petri-net variants. The checker
explores reachable markings, records traces that reach the final marking, and
omits invisible transitions from those traces. It also reports whether the
search finished exhaustively or encountered a resource limit.

The check has explicit resource limits:
\begin{itemize}
  \item at most 5{,}000 explored states;
  \item at most 10{,}000 complete traces;
  \item at most 32 visible activities per trace.
\end{itemize}
The returned trace set must be nonempty and identical for every variant. If all
searches finish, the certificate status is \texttt{exact}; if a bound truncates
at least one search, the compared finite result receives status
\texttt{bounded}. The current configuration sets
\texttt{exact\_equivalence\_only\_for\_training} to false, so a bounded family
is eligible for any split. In the persisted data used here, however, every
controlled family received an \texttt{exact} certificate; there are no bounded
rows.

Ordinary loop families are handled separately because their language is
unbounded. Their second Petri net is an isomorphic clone of the PM4Py tree
conversion, so equivalence is certified structurally by the renaming
transformation. All ordinary-tree rows in the persisted data therefore store
the status \texttt{isomorphic}. Each accepted row includes the certificate
status, semantics, checked variant names, and reference-language size in its
metadata.

The generator also checks simple structural health conditions. In particular,
the final marking must be reachable and there must be no dead transitions. A
family that fails these checks is rejected.

\subsection{Step 5: Create the Generated Log View}
\label{app:create-log-view}

After equivalence is checked or certified, the generator creates a trace pool.
For the controlled families in the current data, the pool is the sorted exact
visible language returned by the completed checker. For ordinary loop
families, PM4Py performs $\max(256,4n)$ process-tree playouts for a requested
log size $n$; the generator deduplicates and sorts the observed traces. Thus,
for the current $n=128$ run, each ordinary trace pool is formed from 512
playout traces. The current run creates one log view per behavior family and
stores 128 traces in that log view.

The default log sampling mode is \texttt{uniform\_variants}. In this mode, the
generator cycles through the trace pool until it has 128 traces, then
shuffles their order with a deterministic seed. If a behavior has only two trace
variants, each variant appears several times. If it has more variants, the first
128 positions of the repeated cycle are used. This keeps the generated logs
simple and reproducible.

The configuration also defines other sampling modes, such as sparse,
long-tail, incomplete, resampled, and noisy modes. They are part of the codebase
for other experiments, but the current persisted split uses one clean
uniform-variants log view per behavior family.

Family construction, quota order, and log-view shuffling use stable seeds
derived from the global seed, split, family index, and pipeline stage. Ordinary
loop trace pools additionally depend on PM4Py process-tree playout, so exact
playout results can depend on the installed PM4Py version even when the Python
seed and generator configuration are unchanged.

\subsection{Step 6: Flatten Families into Training Rows}
\label{app:flatten-rows}

The final step turns each behavior family into rows. With the current
configuration, each family has one generated log view and two Petri-net
variants, so it becomes exactly two rows:
\[
  (\text{tree}, \text{same log view}, \text{Petri variant 0}, e)
  \quad\text{and}\quad
  (\text{tree}, \text{same log view}, \text{Petri variant 1}, e).
\]
Both rows have the same equivalence identifier $e$. The process tree and trace
list are identical in the two rows. The Petri-net graph is different.

Rows are written to:
\begin{itemize}
  \item \texttt{data/training/samples.jsonl};
  \item \texttt{data/validation/samples.jsonl};
  \item \texttt{data/test/samples.jsonl}.
\end{itemize}
The generator also writes \texttt{data/metadata.json}, which records the seed,
the full synthetic configuration, split sizes, motif counts, representation
counts, average tree and trace statistics, and maximum Petri graph sizes.
For each split it records both planned and actual motif quotas, per-motif
representation-slot counts, any strict-coverage deficits, and whether the
minimum-coverage check passed. During generation, \texttt{sample.py} reports
completed flattened rows with a per-split progress bar on standard error;
\texttt{--quiet} disables this display without changing the generated data.

\subsection{How the Row Is Used During Training}
\label{app:training-use}

At training time, each row is converted into tensors:
\begin{itemize}
  \item the process tree becomes a prefix token sequence with explicit arity
        tokens;
  \item the generated traces become padded activity-token sequences;
  \item the Petri net becomes node-type indices, transition-label ids, marking
        features, and two adjacency matrices.
\end{itemize}
The batch collator also creates a positive-pair mask from the equivalence
identifiers. If two rows have the same equivalence identifier, they are treated
as matching views of the same behavior. Training batches can also apply a
consistent random renaming of activity labels within a family. This preserves
the behavior while reducing dependence on the literal label names.

In short, the data-generation pipeline makes the supervision explicit. The
model sees a tree, a finite generated log, and a Petri-net graph; it also sees
which rows describe the same behavior. The generated families therefore train
the model to align artifact types by behavior rather than by node names,
transition counts, or one particular Petri-net structure.
